\documentclass[crop,tikz,11pt]{standalone}
\usepackage{geometricfigures}
\usepackage{amsmath,amssymb}
\usepackage{pgfplots}
\pgfplotsset{compat=1.18}
\usetikzlibrary{arrows.meta,decorations.markings,patterns,intersections}
\usepgfplotslibrary{fillbetween}

\definecolor{lib}{HTML}{3B75AF}
\definecolor{rot}{HTML}{C1701E}
\definecolor{sep}{HTML}{B03030}
\definecolor{net}{HTML}{4E9A4E}

% The pendulum potential cos(theta) read as mu B along a field line, in units in which
% H = v_par^2/2 + cos(theta): the field is strongest on the high-field side theta = 0, 2pi and
% weakest on the low-field side theta = pi. A particle of energy H moves where cos(theta) <= H.
\begin{document}
\begin{tikzpicture}
\begin{axis}[width=11cm,height=5.4cm,xmin=0,xmax=6.2832,ymin=-1.25,ymax=2.0,
  xtick={0,3.1416,6.2832},xticklabels={HFS,LFS,HFS},
  ytick={-1,0,1},
  xlabel={poloidal angle along the field line, $\theta$},ylabel={$\mu B$, energy},
  axis lines=left,clip=false,every axis plot/.append style={samples=200}]
  \addplot[fg,thick,smooth,domain=0:6.2832]{cos(deg(x))};
  % trapped: H = -0.5 between the bounce points arccos(-0.5) and 2pi - arccos(-0.5)
  \addplot[lib,very thick,domain=2.0944:4.1888]{-0.5};
  \fill[lib] (axis cs:2.0944,-0.5) circle (1.6pt) (axis cs:4.1888,-0.5) circle (1.6pt);
  \node[lib,above] at (axis cs:3.1416,-0.5) {\small trapped: bounces};
  % the separatrix touches the maxima, where the particle takes infinitely long to arrive
  \addplot[sep,very thick,dashed,domain=0:6.2832]{1};
  \node[sep,above] at (axis cs:3.1416,1) {\small trapped--passing boundary, $H=1$};
  % passing
  \addplot[rot,very thick,domain=0:6.2832]{1.6};
  \node[rot,above] at (axis cs:3.1416,1.6) {\small passing: circulates};
  \fill[sep] (axis cs:0,1) circle (1.8pt) (axis cs:6.2832,1) circle (1.8pt);
  \node[sep,below right] at (axis cs:0,0.75) {\small $\theta=0$};
  \node[sep,below left] at (axis cs:6.2832,0.75) {\small $\theta=2\pi$};
\end{axis}
\end{tikzpicture}
\end{document}
